Success Rate, First-Attempt Success Rate, Iteration-Weighted Success Rate, Retry Gap --- and the iteration profile that distinguishes genuine correction from stochastic resampling.

This family answers \emph{"how often, how cheaply, and how honestly does the model succeed?"} All four metrics live in the \code{overall}/\code{by\_domain} tables and share one design principle: a success that needed many retries is worth less than a success on the first try, because under the retrieval hypothesis those retries are independent samples, not corrections.

\begin{bmkeybox}{Real values from the bundled run}
gpt\_oss\_120b is the clear leader: SR

0.524, FASR

0.333, IWSR

0.371, Retry

Gap

0.190. qwen\_3\_5\_122b shows the diagnostic pattern that motivates the whole family --- SR

0.500 but FASR

0.000 and Retry

Gap

0.500: every success arrived only after retries, the signature of a Stochastic Searcher.
\end{bmkeybox}

\begin{bmmcard}{METRIC \(\cdot\) SR: Success Rate}
\emph{overall[" Success\_Rate"] = mean(Valid)}
Fraction of instances the model ever solves, across all allowed iterations. The headline number --- and the one the framework argues is least trustworthy on its own.
\begin{description}
\item[How it is computed] Boolean \code{Valid}(from \code{scored["solved"]}) averaged per group. This is \emph{validity@k}: success anywhere in the repair budget counts.

\begin{lstlisting}
Success_Rate=("Valid", "mean")
\end{lstlisting}
\item[Range \& meaningful bands] Bands used by the profile thresholds: \(\approx 0\) \(\rightarrow\) No Grounding \(\cdot\) \(<0.2\) \(\rightarrow\) Understander/Vocabulary \(\cdot\) \(>0.5\) capable \(\cdot\) \(>0.7\) strong. Always read against FASR; SR alone cannot distinguish retrieval from reasoning.
\item[In the plots] \raggedright
\code{success\_rate\_heatmap} (model\(\times\)domain),\allowbreak{}
the SR bars in \code{sr\_vs\_fasr} and
\code{sr\_fasr\_iwsr\_by\_model},\allowbreak{}
and it drives the colour in \code{failure\_mode\_taxonomy} and
\code{fasr\_iwsr\_scatter}.
\item[Rationale] SR is the necessary baseline but the framework's thesis is that it is \emph{insufficient}: it conflates first-try success with retry-dependent success and says nothing about \emph{why} a failure happened. Every other metric exists to decompose it.
\end{description}
\end{bmmcard}

\begin{bmmcard}{METRIC \(\cdot\) FASR: First-Attempt Success Rate}
\emph{overall[" FASR"] = mean(\_iter1)}
Fraction solved on iteration 1 with no feedback at all. The framework treats this as the single honest signal of planning ability, because it is the only one that cannot be inflated by retries.
\begin{description}
\item[How it is computed] A private boolean column is set at row level and averaged:

\begin{lstlisting}
df["_iter1"] = (df["Valid"] == True) & (df["Iterations"] == 1)
FASR=("_iter1", "mean")
\end{lstlisting}
\item[Range \& meaningful bands] \(\mathrm{FASR}>0.50\) with low Retry Gap \(\rightarrow\) genuine zero-shot capability. FASR \(\approx\) 0 while SR is high \(\rightarrow\) the model only ever succeeds by resampling (the entire qwen / nemotron pattern in the bundled run).
\item[In the plots] \raggedright
\code{fasr\_by\_model\_domain},\allowbreak{}
\code{fasr\_by\_difficulty},\allowbreak{}
the FASR bar in \code{sr\_vs\_fasr}/\allowbreak{}\code{sr\_fasr\_iwsr\_by\_model},\allowbreak{}
both axes of \code{fasr\_iwsr\_scatter},\allowbreak{}
and the y-axis of \code{failure\_mode\_taxonomy} bubble size.
\item[Rationale] Under the retrieval hypothesis (Kambhampati 2024) and the irrelevant-feedback finding (Stechly 2023), only the un-retried attempt reflects what the model can do in one shot. FASR carries the highest weight in the composite score (0.25) for exactly this reason.
\end{description}
\end{bmmcard}

\begin{bmmcard}{METRIC \(\cdot\) IWSR: Iteration-Weighted Success Rate}
\emph{contrib = (1/Iterations) if Valid else 0}
Success discounted by how many attempts it took. Rewards models that converge fast; penalises models that only succeed deep into the retry budget.
\begin{description}
\item[How it is computed] Each solved instance contributes \code{1/k} where k is the iteration count; unsolved contributes 0. The mean of these contributions is IWSR.

\begin{lstlisting}
df["_iwsr_contrib"] = (1.0/row["Iterations"]) if row["Valid"] and row["Iterations"]>0 else 0.0
IWSR=("_iwsr_contrib", "mean")
\end{lstlisting}
\item[Range \& meaningful bands] IWSR is bounded above by SR and below by FASR (a first-try win contributes the full 1.0, matching FASR; a 5th-try win contributes only 0.2). IWSR \(\approx\) SR \(\rightarrow\) fast repairs (Efficient Corrector). IWSR \(\ll\) SR \(\rightarrow\) slow, late convergence.
\item[In the plots] \code{iwsr\_by\_model\_domain}, the IWSR bar in \code{sr\_fasr\_iwsr\_by\_model}, and the y-axis of \code{fasr\_iwsr\_scatter}.
\item[Rationale] The \code{1/k} weighting is the numerical encoding of Stechly's argument: if verification isn't easier than generation for a retrieval system, late successes deserve no credit for "search". IWSR is the metric that turns the retry budget into a cost.
\end{description}
\end{bmmcard}

\begin{bmmcard}{METRIC \(\cdot\) RG: Retry Gap}
\emph{Retry\_Gap = Success\_Rate - FASR}
How much of the reported success depends on retries. A pure diagnostic, not a capability --- high is a warning, not a strength.
\begin{description}
\item[How it is computed] A single subtraction on the \code{overall}/\code{by\_domain} tables, then sorted ascending in the \code{retry\_gap} table.

\begin{lstlisting}
overall["Retry_Gap"] = overall["Success_Rate"] - overall["FASR"]
\end{lstlisting}
\item[Range \& meaningful bands] \(\mathrm{RG}<0.10\) \(\rightarrow\) near-zero-shot. 0.10--0.30 \(\rightarrow\) Efficient Corrector territory. \(\mathrm{RG}>0.30\) \(\rightarrow\) Stochastic Searcher: success is mostly resampling. In the bundled run qwen\_3\_5 hits 0.500.
\item[In the plots] \code{retry\_gap\_by\_model}(diverging horizontal bars, red = positive gap). RG is also a ranked metric (direction = min) in the within-domain ranking.
\item[Rationale] RG operationalises the difference between "solved" and "solved on its own". It is the cleanest single read on whether a model's SR is trustworthy.
\end{description}
\end{bmmcard}

\subsubsection{The iteration profile --- the test that separates correction from resampling}

RG tells you retries happened; the iteration profile tells you \emph{whether they helped}.\code{compute\_iteration\_profile()} computes, for each attempt slot k, the conditional success probability among instances that needed exactly k attempts:

\begin{lstlisting}
for k in range(1, max_iter+1):
    reached = df[df["Iterations"] >= k]
    exact   = reached[reached["Iterations"] == k]
    p_valid_given_reached = exact["Valid"].mean()
\end{lstlisting}

A \textbf{rising} curve means later attempts succeed more often --- the model is using feedback (Efficient Corrector). A \textbf{flat} curve means success probability is independent of attempt number --- each retry is an independent draw (Stochastic Searcher).\code{profile\_is\_flat()} formalises "flat" as max-min \(\le\) 0.10 over the finite points, and feeds the Stochastic-Searcher profile condition. This is the direct empirical test of Stechly's finding (iii) in your own data.

\begin{bmcavbox}{Reading caveat}
With few instances per (model, k) slot the curve is noisy and

\code{profile\_is\_flat}

can return None (inconclusive) when fewer than two finite points exist. Treat the profile as suggestive, not decisive, on small task sets --- the bundled run's 6-instance subsets are exactly this regime.
\end{bmcavbox}

\subsubsection{Reading the efficiency plots}

The three success metrics are only diagnostic in combination. The guarantee \code{FASR \(\le\) IWSR \(\le\) SR} holds by construction, and the two gaps between them carry the meaning: the \emph{SR--IWSR} gap measures how late successes arrive, the \emph{IWSR--FASR} gap measures how much retries contributed at all.

\paragraph{FASR by model \(\times\) domain.}

Two readings live in one chart.\textbf{Per model (compare a model's bars across domains):} a systematic gap between one domain group and another suggests memorised patterns specific to that group rather than generalised planning; a small gap across all domains points to genuine zero-shot transfer. A uniformly low FASR means near-zero zero-shot ability --- go to SR-vs-FASR to check whether success is being manufactured by resampling.\textbf{Per domain (compare one domain across models):} a domain with consistently high FASR across every model is simply an easy domain --- the domain, not the model, is driving the score.

\paragraph{SR vs FASR --- the retry gap.}

\begin{table}[h]\small
\centering
\begin{tabularx}{\linewidth}{@{}l Y@{}}
\toprule
Gap pattern & Meaning\\
\midrule
SR \(\approx\) FASR, both high & Genuine zero-shot capability. Retries add almost nothing.\\
\addlinespace
\(\mathrm{SR}\) moderately \(>\) \(\mathrm{FASR}\) & Partial capability, partly scaffolded by the retry mechanism.\\
\addlinespace
SR \(\gg\) FASR (large gap) & Stochastic search. Per Stechly et al. (2023), retries are re-sampling, not correction --- SR alone misleads.\\
\addlinespace
Both near 0 & No planning capability; the retry budget is irrelevant.\\
\bottomrule
\end{tabularx}
\end{table}

\paragraph{FASR by difficulty bin --- genuine reasoning vs memorisation.}

An easy\(\rightarrow\)hard sweep separates reasoning from training-frequency exploitation. A reasoner degrades \emph{gradually and smoothly} as steps, objects and constraint interactions grow --- the decline is a property of the task. A memoriser stays high on easy (short, common-looking plans) and then falls off a \emph{cliff} at the point where the required plans become long and rare, i.e. fall outside its training distribution. The cliff is the boundary of the distribution, not of a reasoning process.

\begin{table}[h]\small
\centering
\begin{tabularx}{\linewidth}{@{}Y I P@{}}
\toprule
FASR pattern (easy\(\rightarrow\)hard) & Interpretation & Reference\\
\midrule
Gradual decline & Generalisation --- graceful degradation under complexity; some genuine reasoning that scales. & Expected for capable reasoners\\
\addlinespace
Sharp cliff to \textasciitilde{} 0 & Memorisation --- easy instances exploit training frequency; hard instances fall outside the distribution. & McCoy et al. 2024\\
\addlinespace
Flat and high & Genuine capability across all difficulties. Strongest possible result. & Ideal case\\
\addlinespace
Flat and low & No capability at any difficulty --- or Difficulty='unknown' throughout (check data). & Valmeekam et al. 2023\\
\addlinespace
Non-monotonic / irregular & Difficulty bins don't track real plan complexity. Recalibrate using optimal-plan-length quantiles / statistical rarity. & Data-quality issue\\
\bottomrule
\end{tabularx}
\end{table}

\paragraph{IWSR --- where in the run the success happened.}

The closer IWSR sits to FASR, the more successes landed on attempt 1; the closer it sits to SR, the faster the model converges when it does eventually succeed. The two gaps read together:

\begin{table}[h]\small
\centering
\setlength{\tabcolsep}{4pt}
\begin{tabularx}{\linewidth}{@{}C C I P@{}}
\toprule
SR--IWSR & IWSR--FASR & Interpretation & Profile\\
\midrule
small & small & All three nearly equal; most success on attempt 1, retries add nothing. & Genuine zero-shot planner\\
\addlinespace
small & large & IWSR \(\approx\) SR \(\gg\) FASR --- rarely first-try, but converges quickly on attempt 2--3. & Efficient corrector\\
\addlinespace
large & small & FASR \(\approx\) IWSR \(\ll\) SR --- most success arrives very late, at attempt 4--5, and is heavily penalised by IWSR. & Stochastic / late luck\\
\addlinespace
large & large & All three spread apart; SR is inflated by retries, while FASR reveals that true zero-shot performance is low. & Fully retry-dependent\\
\bottomrule
\end{tabularx}
\end{table}

\paragraph{P(Valid | reached attempt k).}

For each k the curve takes plans that actually \emph{reached} attempt k (\code{Iterations \(\ge\) k}) and reports the fraction resolving at exactly k --- a conditional, not cumulative, probability. Its shape is the direct empirical test of the Stochastic-Searcher hypothesis.

\begin{table}[h]\small
\centering
\begin{tabularx}{\linewidth}{@{}Y I P@{}}
\toprule
Line shape & Interpretation & Reference\\
\midrule
Declining from k=1 & Easiest problems resolve first; harder ones take more attempts --- difficulty-sensitive correction. & Expected for a capable reasoner\\
\addlinespace
Flat across k & Each retry is statistically independent --- re-sampling, not correcting. Feedback content is irrelevant. & Stechly et al. 2023 (\S{}4 Results)\\
\addlinespace
Increasing with k & Model improves with each retry --- genuine in-context learning. Very rare; check whether feedback leaks the answer. & Would contradict Stechly et al.\\
\addlinespace
Low-flat near 0 & No valid plans at any attempt; the iteration budget is wasted compute. & Valmeekam et al. 2023 (\S{}5.2)\\
\bottomrule
\end{tabularx}
\end{table}

\paragraph{Retry-gap bar \& the FASR--IWSR scatter.}

The diverging retry-gap bar colours each model by \code{SR - FASR}\allowbreak{}(tomato for the usual positive gap; a steelblue/negative bar would be a data anomaly, since SR \(\ge\) FASR by construction). Models at the top --- lowest gap --- are the most independent of the retry scaffold. The FASR-vs-IWSR scatter then encodes three axes at once:\textbf{x = FASR}(zero-shot skill),\textbf{y = IWSR}(retry efficiency),\textbf{dot size+colour = SR}(deep-red small = SR 0 \(\rightarrow\) deep-green large = SR 1). No point can fall below the FASR=IWSR diagonal.

\begin{table}[h]\small
\centering
\begin{tabularx}{\linewidth}{@{}l l Y@{}}
\toprule
Position & Dot & Interpretation\\
\midrule
Top-right, on/near diagonal & large green & Genuine planner --- high FASR, IWSR\(\approx\)FASR, high SR; success is first-attempt.\\
\addlinespace
Top-left, well above diagonal & medium yellow & Efficient corrector --- low FASR lifted by early-retry successes.\\
\addlinespace
Bottom-left & small red & No capability --- FASR and IWSR near zero, SR low.\\
\addlinespace
Top-left, large dot, low FASR & large green/yellow & Dangerous case --- high SR/IWSR mask retry dependence; SR misleads without FASR.\\
\bottomrule
\end{tabularx}
\end{table}
